% Mathématiques 6e — cours V3
% Source LaTeX rendue côté Astro/KaTeX.

\section{Objectifs}
\begin{itemize}
\item Reconnaître une situation de proportionnalité.
\item Utiliser un coefficient ou le passage à l'unité.
\item Exploiter une échelle.
\item Résoudre des problèmes de pourcentages simples.
\end{itemize}

\section{1. Situation de proportionnalité}
Deux grandeurs sont proportionnelles lorsqu'on passe de l'une à l'autre en multipliant toujours par le même nombre.

Si 3 kg de pommes coûtent 7,50 €, alors 1 kg coûte
\[
7{,}50\div3=2{,}50\text{ €}.
\]
Pour 5 kg :
\[
5\times2{,}50=12{,}50\text{ €}.
\]

\section{2. Coefficient de proportionnalité}
Dans une situation proportionnelle :
\[
y=kx.
\]
Le nombre $k$ est le coefficient de proportionnalité.

En 6e, on privilégie le sens des grandeurs, le passage à l'unité, la linéarité et les coefficients simples.

\section{3. Échelle}
Une échelle $1:50\,000$ signifie que 1 unité sur la carte correspond à 50 000 mêmes unités dans la réalité.

Ainsi, 1 cm sur la carte représente 50 000 cm, soit 500 m.

\section{4. Pourcentages}
Calculer 20 \% de 150 peut se faire en remarquant que 10 \% vaut 15, donc 20 \% vaut 30.

\section{5. Limite explicite du programme}
\textbf{Le produit en croix n'est pas enseigné en 6e.}

On utilise des raisonnements accessibles : retour à l'unité, multiplication/division, propriétés de linéarité ou fractions simples.

\section{Erreurs fréquentes}
\begin{itemize}
\item Supposer qu'une situation est proportionnelle sans le vérifier.
\item Mélanger les unités dans un problème d'échelle.
\item Appliquer mécaniquement une technique de produit en croix.
\end{itemize}

\section{À retenir}
En proportionnalité, un même coefficient relie les deux grandeurs. En 6e, on raisonne d'abord sur le sens, l'unité et les facteurs multiplicatifs.
